\documentclass[11pt]{article}
\usepackage[T1]{fontenc}
\usepackage{lmodern,amsmath,amssymb,amsthm,mathtools,longtable,array}
\usepackage[margin=1.05in]{geometry}
\usepackage[hidelinks]{hyperref}
\newtheorem{theorem}{Theorem}[section]
\newtheorem{lemma}[theorem]{Lemma}
\newtheorem{proposition}[theorem]{Proposition}
\newtheorem{corollary}[theorem]{Corollary}
\theoremstyle{remark}\newtheorem{remark}[theorem]{Remark}
\newcommand{\R}{\mathbb R}\newcommand{\N}{\mathbb N}
\newcommand{\D}{\mathcal D}\newcommand{\E}{\mathcal E}
\newcommand{\norm}[1]{\left\lVert#1\right\rVert}
\newcommand{\code}[1]{{\small\nolinkurl{#1}}}
\title{Signed multichannel reconstruction\protect\\
on a factorial analytic cylinder domain}
\author{Project Stardust}
\date{Reviewed editorial draft, October 29, 2026}
\begin{document}
\maketitle
\begin{abstract}
For the signed heat kernel of even order $p=2m$, we construct a
multichannel reconstruction formula on a specified Banach completion
of dyadic polynomial cylinders. The coefficient fields are genuine
Bochner $L^1$ limits of native predictable Taylor coefficients and
have degree-independent bounds on the original observable norm.
The centered observable pairing is the sum of $p-1$ bounded signed
channels in the continuous test dual. The proof begins with an exact
finite signed telescope, proves quantitative increment and
adjacent-refinement estimates, and only then passes to the observable
completion. A joint diagonal realizes every completed observable
through actual signed polynomial sums. The pairing is indefinite;
no positive path-space measure or stochastic isometry is asserted.
We distinguish canonical coefficient uniqueness from uniqueness of channel synthesis. Exact variation/path-law, quartic-error, deterministic $L^2$-input and terminal-channel obstructions delimit stronger interpretations. All principal results are supported by a frozen
Lean development. Novelty relative to earlier higher-order signed
calculus remains unresolved.
\end{abstract}

\section{Scope and normalization}
Fix $m\in\N$, $m>0$, a physical horizon $T>0$, and put
\[
 p=2m,\qquad c=(-1)^{m+1}.
\]
The fundamental signed kernel $k_\tau$, $\tau>0$, has Fourier
transform $e^{-\tau\xi^p}$ and generator $c\partial_x^p$.
We normalize the Fourier transform by
\[
 k_\tau(x)=\frac1{2\pi}\int_\R e^{-\tau\xi^p}e^{-i\xi x}\,d\xi.
\]
Changing the Fourier sign gives the same even real kernel.
For each finite partition, independent increments mean the product
of these signed densities; they do not mean independent random
variables under an assumed positive path law.

The analytic domain $\D$ will be defined explicitly in Section~\ref{sec:domain}.
It has a dense injective polynomial core, a bounded signed mean
$\E_T:\D\to\R$, and a bounded symmetric pairing
\[
 J_T:\D\longrightarrow\D^*,\qquad
 J_T(P)(G)=\E_T(PG)
\]
on polynomial cylinders. Products in this last identity are finite
polynomial products; bounded multiplication on the whole completion
is not being assumed. Let $\lambda$ denote Lebesgue measure on
$[0,1]$, so time in the coefficient fields is normalized time and
the physical horizon occurs in all kernel parameters and channel
constants. Write $\mathcal F=L^1(\lambda;\D^*)$ for Bochner $L^1$.

\begin{theorem}\label{thm:main}
There are canonical bounded linear maps $C_q:\D\to\mathcal F$,
$q\geq0$, satisfying
\begin{equation}\label{eq:Cbound}
 \norm{C_q u}_{\mathcal F}\leq
 \frac{\norm{J_T}e^T}{q!}\norm u_{\D}.
\end{equation}
For $1\leq q\leq p-1$, define
\begin{equation}\label{eq:channel}
 (I_q A)(v)=cpT(p-1)_{q-1}
 \int_0^1 A(s)(D_s^{p-q}v)\,ds,
\end{equation}
where $(r)_k=r(r-1)\cdots(r-k+1)$ and $D_s$ is the bounded cutoff
derivative defined below. Then every $u\in\D$ satisfies
\begin{equation}\label{eq:main}
 J_T\big(u-\E_T(u)1\big)
 =\sum_{q=1}^{p-1} I_q(C_q u)\qquad\text{in }\D^*.
\end{equation}
On every fixed polynomial cylinder, $C_q$ is the norm limit in
$\mathcal F$ of its native predictable Taylor coefficient fields.
Its values lie in the closed span of actual adapted polynomial
step fields. Every completed observable admits a simultaneous
polynomial/mesh approximation of all surviving fields and the
right side of \eqref{eq:main}.
\end{theorem}

The equality is in the test dual, not an assertion of equality in
$\D$ obtained by inverting $J_T$. Injectivity of $J_T$ on the entire
completion is not supplied. The theorem is therefore a signed
Clark--Ocone-type reconstruction with its topology specified.

Higher-order signed heat calculus has substantial antecedents,
including Hochberg's work~\cite{hochberg}. Complex random-walk
constructions give another generalized calculus~\cite{bcm}, while
time-changed pseudoprocesses are studied in~\cite{co}. The present
paper does not claim the first higher-order It\^o calculus. Its
precise subject is the simultaneous native coefficient limit and
bounded observable completion for these actual signed cylinder
laws. Priority for this package remains unresolved.

\section{Actual finite signed laws and polynomial conditioning}
The function $e^{-\tau\xi^p}$ is Schwartz. Fourier inversion shows
that $k_\tau$ is real, even and Schwartz, has integral one, and
satisfies $k_\tau*k_\sigma=k_{\tau+\sigma}$. Thus polynomial
functions are integrable against its total variation. Its moments
are
\begin{equation}\label{eq:moments}
 \mu_r(\tau)=\int x^r k_\tau(x)\,dx
 =\begin{cases}(pk)!(c\tau)^k/k!,&r=pk,\\0,&p\nmid r.
 \end{cases}
\end{equation}
For $r=0$ this gives one. To obtain the formula, differentiate the
Fourier transform at zero; Schwartz decay justifies differentiation
under the integral. The derivative of $e^{-\tau\xi^p}$ vanishes
unless the order is a multiple of $p$, and the resulting powers
of $i$ give $c^k$. Equivalently, for every polynomial $f$,
\begin{equation}\label{eq:heat}
 \int f(x)k_\tau(x)\,dx
 =\big(e^{c\tau\partial^p}f\big)(0).
\end{equation}
The exponential series terminates on each polynomial.

An integration-by-parts identity that will control all power
channels is
\begin{equation}\label{eq:stein}
 \int xf(x)k_\tau(x)\,dx
 =cp\tau\int f^{(p-1)}(x)k_\tau(x)\,dx.
\end{equation}
It follows either from the moments \eqref{eq:moments}, by testing
monomials, or from differentiation of the Fourier transform.
All boundary terms vanish by Schwartz decay.

On the dyadic grid with $N=2^n$ cells set $h=1/N$, $\tau=Th$,
and use increment variables $x_0,\ldots,x_{N-1}$. The signed
expectation $\E_{T,n}$ integrates against
$\prod_j k_\tau(x_j)\,dx_j$. The position variables are
$y_i=x_0+\cdots+x_i$. Refinement replaces each coarse increment
by the sum of two fine increments. Convolution implies consistency
of these actual signed expectations.

For $0\leq j\leq N$, let $\Pi_jP$ integrate increments with
indices $i\geq j$, leaving $x_0,\ldots,x_{j-1}$ untouched.
It is a polynomial in the retained increments. In particular
$\Pi_0P=\E_{T,n}P$, $\Pi_NP=P$, and $\Pi_j\Pi_k=\Pi_{\min(j,k)}$.
These formulas follow from finite signed Fubini, with absolute
integrability of polynomials against the product variation.
The notation is polynomial conditioning; no general conditional
expectation theorem for an indefinite path law is invoked.

\section{The common factorial domain}\label{sec:domain}
Let $S$ be the union of the positive dyadic grid times in $(0,1]$.
A cylinder is a polynomial in finitely many formal position
coordinates $Y_s$, $s\in S$. The direct-limit construction using
increment refinement gives exactly this polynomial algebra: old
positions are retained at their original times, and the cumulative
coordinate change is invertible on each finite grid.

For a finite multiindex $\alpha:S\to\N$, let
$|\alpha|=\sum_s\alpha_s$. For a polynomial
$P=\sum_\alpha a_\alpha Y^\alpha$, define
\begin{equation}\label{eq:norm}
 \norm P_{\D}=\sum_\alpha |\alpha|!\,|a_\alpha|.
\end{equation}
This is the factorial of total degree, not the product of coordinate
factorials. The coefficients here are in \emph{position} coordinates,
after converting a given increment polynomial by the inverse
cumulative substitution. Refinement is an isometry because it only
retains the same position labels. It also does not increase total
degree. The norm is nondegenerate because polynomial coefficients
are unique.

Define $\D$ as the completion of this normed polynomial space.
More concretely, it is the closed span of the vectors
$(|\alpha|!a_\alpha)_\alpha$ in $\ell^1$ indexed by finitely
supported time multiindices. The Lean construction uses all times
in $(0,1]$ as an ambient index and takes the closed span of the
dyadic core. This yields the same space just described. It is a
specified analytic coefficient completion, not a space of all
bounded smooth path observables.

\begin{lemma}\label{lem:operators}
For any direction $v:S\to\R$ with $|v_s|\leq1$, the polynomial
directional derivative $\partial_v=\sum_s v_s\partial_{Y_s}$
extends to $\D$ with norm at most one. Position substitutions that
send each variable to another variable or zero, without scaling
by a coefficient of modulus greater than one, are contractions.
The finite signed projection $\Pi_j$ has norm at most $e^T$ on
each cylinder norm.
\end{lemma}
\begin{proof}
For a monomial of degree $d>0$, the ratio of the weighted norm of
its directional derivative to $d!$ is at most
$\sum_s\alpha_s/d=1$. The degree-zero derivative vanishes.
Summation proves the derivative bound. Substitutions preserve
total degree or lower it; merging equal monomials can only reduce
the sum of absolute coefficients.

In position coordinates, differentiation in the increment $x_j$
is the direction which is one on positions after that increment
and zero before it. Its norm is therefore at most one.
Integrating one increment is heat evolution
$e^{c\tau\partial_{x_j}^p}$ followed by setting that increment
to zero. Setting it to zero merges later positions with their
predecessors, so is contractive. The heat operator has norm at
most $e^\tau$ by its exponential series. Composing future heat
operators has norm at most $e^{\sum\tau}\leq e^T$, and the
remaining clamping substitution is contractive.
\end{proof}

Taking the empty-past projection gives a bounded signed mean
$\E_T$ on $\D$, with $\norm{\E_T}\leq e^T$ and $\E_T1=1$.
The cutoff derivative in \eqref{eq:channel} is
\begin{equation}\label{eq:Ds}
 D_s=\partial_{v(s)},\qquad v(s)_t=1_{\{s\leq t\}}.
\end{equation}
It is contractive, as are all its powers. Applied to a fixed
polynomial, it is a finite step function of $s$. Approximation
by core polynomials shows $s\mapsto D_s^r u$ is strongly measurable
for every $u\in\D$: the applied fields converge pointwise, even
uniformly in the norm bound. This gives a bounded map into
$L^\infty(\lambda;\D)$.

\section{The product pairing and its signed character}
Ordinary multiplication does not have norm at most the product
of the norms in \eqref{eq:norm}. Instead define on polynomials
\[
 B(P,Q)(Y)=P(Y/2)Q(Y/2).
\]
For monomials of degrees $d,e$, its weighted norm ratio is
$2^{-d-e}(d+e)!/(d!e!)\leq1$. Summing absolute coefficients
gives $\norm{B(P,Q)}\leq\norm P\norm Q$, so $B$ extends
as a bounded symmetric bilinear map $\D\times\D\to\D$.
Set
\begin{equation}\label{eq:pairing}
 J_T(u)(v)=\E_{2^pT}\big(B(u,v)\big).
\end{equation}
Then $\norm{J_T}\leq e^{2^pT}$. Homogeneity of the moments shows
that half-dilation at horizon $2^pT$ restores horizon $T$:
$\mu_r(2^p\tau)2^{-r}=\mu_r(\tau)$.
Applying this to each increment monomial proves
\begin{equation}\label{eq:actualpair}
 J_T(P)(G)=\E_T(PG)
\end{equation}
for every pair of cylinders, and $J_T(u)(1)=\E_Tu$ by continuity.
This construction supplies the bounded product \emph{pairing}
without assuming bounded ordinary multiplication.

For $p=4$ one has $c=-1$, so \eqref{eq:moments} gives
\[
 J_T(Y_1^2)(Y_1^2)=\E_T(Y_1^4)=-24T<0.
\]
Thus the pairing is indefinite. In particular the notation
$\E_T$ cannot be used to justify positive Cauchy--Schwarz,
martingale isometries, or Hilbert orthogonal projections.

For later estimates it is useful to name
\[
 W_j(H)(s)=1_{(jh,(j+1)h]}(s)J_T(H).
\]
This is an actual element of $\mathcal F$, with
$\norm{W_j(H)}_{\mathcal F}=h\norm{J_T(H)}$.
An adapted step field has this form summed over $j$, with $H_j$
a polynomial depending only on increments before cell $j$.
Let $\mathcal A_T$ be the closed linear span of all such fields
in $\mathcal F$. This definition is the exact adaptedness space
used below; it does not presuppose a probability filtration.

\section{The bounded channel operators}
For $A\in\mathcal F$ and $v\in\D$, the integrand
$A(s)(D_s^r v)$ is integrable because its absolute value is at
most $\norm{A(s)}\norm v$. Strong measurability follows from
the bounded bilinear evaluation map and the applied-field
measurability proved above. Define
\[
 Q_r(A)(v)=\int_0^1 A(s)(D_s^r v)\,ds.
\]
Then $Q_r:\mathcal F\to\D^*$ is bounded with norm at most one.
In addition to the surviving channels in \eqref{eq:channel}, use
the same formula for $q=p$. In that case
\begin{equation}\label{eq:driftchannel}
 I_p=cp!T Q_0.
\end{equation}
It is a drift channel, which will cancel exactly after centering.
For $1\leq q\leq p$,
\begin{equation}\label{eq:Ibound}
 \norm{I_q A}\leq pT(p-1)_{q-1}\norm A_{\mathcal F}.
\end{equation}

For recursive calculations put
$K_k=cpT\binom{p-1}{k}Q_{p-1-k}$, $0\leq k<p$.
Thus $I_q=(q-1)!K_{q-1}$.
The factorial distinction is needed: the native Taylor coefficient
contains $1/q!$, whereas conversion of a power channel contributes
$(q-1)!$. Losing either factor changes the reconstruction formula.

\section{Exact finite predictable innovations}
Fix a grid and a polynomial $P$ of total degree at most $d$.
For cell $j$, put $Q_j=\Pi_{j+1}P$, and let $R_j$ set $x_i=0$
for $i\geq j$. Define the native coefficients
\begin{equation}\label{eq:nativejet}
 a_{j,q}(P)=\frac1{q!}R_j\partial_{x_j}^q Q_j.
\end{equation}
They depend only on strict-past increments, vanish for $q>d$,
and Lemma~\ref{lem:operators} gives
\begin{equation}\label{eq:jetbound}
 \norm{a_{j,q}(P)}_{\D}\leq\frac{e^T}{q!}\norm P_{\D}.
\end{equation}
Only the current and past increments survive in $Q_j$.
Its exact Taylor expansion in the current increment therefore is
$Q_j=\sum_{q=0}^d a_{j,q}(P)x_j^q$.
Integrate that increment to obtain
$\Pi_jP=\sum_{q=0}^d a_{j,q}(P)\mu_q(Th)$.
Subtracting proves the cell identity
\begin{equation}\label{eq:cell}
 \Pi_{j+1}P-\Pi_jP
 =\sum_{q=1}^d a_{j,q}(P)\big(x_j^q-\mu_q(Th)\big).
\end{equation}
The zeroth term vanishes because $\mu_0=1$.
Summing over cells telescopes exactly:
\begin{equation}\label{eq:telescope}
 P-\E_TP=\sum_{q=1}^d\sum_{j=0}^{N-1}
 a_{j,q}(P)\big(x_j^q-\mu_q(Th)\big).
\end{equation}
No term has been discarded on a scaling heuristic. Every integral
in this derivation is an actual finite signed polynomial integral.

Refine a fixed base polynomial $P$ to successive grids and denote
its refinement by $P_n$. Its position-coordinate norm is unchanged
and its total degree stays at most $d$. Define the native field
\begin{equation}\label{eq:nativefield}
 A_{n,q}(P)=\sum_j W_j\big(a_{j,q}(P_n)\big).
\end{equation}
Since cell lengths sum to one, \eqref{eq:jetbound} gives
\begin{equation}\label{eq:fieldbound}
 \norm{A_{n,q}(P)}_{\mathcal F}
 \leq\frac{\norm{J_T}e^T}{q!}\norm P_{\D}.
\end{equation}
This estimate is independent of degree and mesh, but does not by
itself prove that the sequence converges. That requires the next
two quantitative arguments.

\section{Signed power estimates in the test dual}\label{sec:powers}
For strict-past $H$ and integer $r\geq0$, signed integration by
parts gives the exact dual identity
\begin{equation}\label{eq:power}
 J_T(Hx_j^{r+1})
 =\sum_{k=0}^{p-1}(r)_k K_k\big(W_j(Hx_j^{r-k})\big).
\end{equation}
Terms with $k>r$ vanish and can be omitted. To check the identity,
test first against a polynomial $G$ on a common finer grid.
The coarse increment is the sum of fine increments in its cell.
Apply \eqref{eq:stein} to each fine increment and use Leibniz's
rule on $x_j^rG$. Derivatives of $H$ vanish because it is strict
past. The $k$ derivatives on $x_j^r$ give $(r)_k x_j^{r-k}$,
the other derivatives give $D_s^{p-1-k}G$, and summing the fine
durations gives exactly the cell integral in $K_k$.
This proves \eqref{eq:power} on every polynomial test. Density and
continuity give the identity in $\D^*$.

Define nonnegative constants recursively by
\begin{equation}\label{eq:powerconstants}
 B_0=1,\qquad B_{r+1}=pT\sum_{k=0}^{\min(r,p-1)}
 \binom{p-1}{k}(r)_k B_{r-k}.
\end{equation}
They depend on $p,T,r$, not on the mesh or $H$.

\begin{lemma}\label{lem:powerbounds}
For a strict-past polynomial $H$ and a cell of normalized length
$h\leq1$,
\begin{align}
 \norm{J_T(Hx_j^r)}&\leq B_r\norm{J_T(H)},\label{eq:uniformpower}\\
 \norm{J_T(Hx_j^q)}&\leq B_qh\norm{J_T(H)}\quad(q\geq1),\label{eq:positivepower}\\
 \norm{J_T(Hx_j^q)}&\leq B_qh^2\norm{J_T(H)}\quad(q>p).\label{eq:highpower}
\end{align}
For $1\leq q\leq p$,
\begin{equation}\label{eq:powererror}
 \norm{J_T(Hx_j^q)-I_q(W_jH)}
 \leq B_qh^2\norm{J_T(H)}.
\end{equation}
\end{lemma}
\begin{proof}
The channel bound for a step is
$\norm{K_k(W_jV)}\leq pT\binom{p-1}{k}h\norm{J_T(V)}$.
Apply it in \eqref{eq:power} and induct on $r+1$, using
\eqref{eq:powerconstants}. This first proves
\eqref{eq:positivepower}; since $h\leq1$, it also proves
\eqref{eq:uniformpower} and closes the induction.
For $q>p$, every surviving index $k\leq p-1$ has $q-1-k\geq1$.
Use \eqref{eq:positivepower} on that residual power in each term
of \eqref{eq:power}. The additional step length supplies $h^2$,
proving \eqref{eq:highpower}.
For $q\leq p$, the last term $k=q-1$ is exactly
$(q-1)!K_{q-1}(W_jH)=I_q(W_jH)$.
Every other term has a positive residual power, so the same
argument bounds their sum by $B_qh^2\norm{J_T(H)}$.
\end{proof}

By \eqref{eq:moments}, $\mu_q(Th)=0$ for $1\leq q<p$ and
$\mu_p(Th)=cp!Th$. For $q>p$, $|\mu_q(Th)|\leq M_qh^2$
with a mesh-independent $M_q$: if $q=pk$, take
$M_q=q!T^k/k!$ and note $k\geq2$; otherwise take zero.
Thus summing over $N=1/h$ cells, under any uniform bound
$\norm{J_T(H_j)}\leq M$, gives the following three limits:
\begin{align}
 J_T\left(\sum_jH_j(x_j^q-\mu_q(Th))\right)
 &\longrightarrow I_q(A),&&1\leq q<p,\label{eq:lowl}\\
 J_T\left(\sum_jH_j(x_j^p-\mu_p(Th))\right)
 &\longrightarrow0,\label{eq:criticall}\\
 J_T\left(\sum_jH_j(x_j^q-\mu_q(Th))\right)
 &\longrightarrow0,&&q>p,\label{eq:highl}
\end{align}
where the first line also assumes $\sum_jW_jH_j\to A$ in
$\mathcal F$. The errors are $O(h)M$ in dual norm.
For the critical line, \eqref{eq:driftchannel} gives
$I_p(W_jH_j)=cp!ThJ_T(H_j)$, which cancels the centering exactly.
This is why the final centered reconstruction has $p-1$ channels.

\section{Causal coefficients and convergence under refinement}
It remains to prove the field convergence needed in
\eqref{eq:lowl}. Introduce the causal comparison coefficient
\begin{equation}\label{eq:causaljet}
 b_{j,q}(P)=\frac1{q!}\Pi_j(\partial_{x_j}^qP),
\end{equation}
and its step field $\widetilde A_{n,q}=\sum_j W_j(b_{j,q}(P_n))$.
The derivative here is taken before the strict-past projection.
It does not in general commute through that projection.

\begin{lemma}\label{lem:heatshift}
For every grid polynomial,
\begin{equation}\label{eq:heatshift}
 b_{j,q}(P)=\frac1{q!}R_j
 e^{cTh\partial_{x_j}^p}\partial_{x_j}^q\Pi_{j+1}P.
\end{equation}
Consequently, for a fixed cylinder under refinement,
\begin{equation}\label{eq:heaterror}
 \norm{\widetilde A_{n,q}-A_{n,q}}_{\mathcal F}
 \leq\frac{\norm{J_T}e^T}{q!}(e^{Th}-1)\norm P_{\D}
 \longrightarrow0.
\end{equation}
\end{lemma}
\begin{proof}
The current coordinate is retained by $\Pi_{j+1}$, so its derivative
commutes with the future integrations. Integrating that one current
coordinate after future conditioning is exactly heat evolution
followed by clamping, by \eqref{eq:heat}. This proves
\eqref{eq:heatshift}. The norm of
$e^{cTh\partial_{x_j}^p}-\mathrm{id}$ is at most $e^{Th}-1$.
Use contractivity of derivatives and clamping, the $e^T$ projection
bound, and the sum of normalized cell lengths equal to one.
\end{proof}

The heat correction proves equality of the eventual native and
causal limits, but does not itself prove that either sequence has
a limit. For that we need control of a moving conditioning cut.

\begin{lemma}\label{lem:cut}
For each degree bound $d$, there is $L_d\geq0$, depending only
on $p,T,d$, such that on every grid and every degree-at-most-$d$
polynomial $P$,
\begin{equation}\label{eq:cut}
 \norm{J_T(\Pi_bP-\Pi_aP)}
 \leq L_d(b-a)h\norm P_{\D},\qquad0\leq a\leq b\leq N.
\end{equation}
\end{lemma}
\begin{proof}
For each $q\geq1$ there is $V_q\geq0$ with
$|\mu_q(Th)|\leq V_qh$: use zero if $p\nmid q$, and
$V_q=q!T^k/k!$ if $q=pk$, since $h^k\leq h$.
The cell identity \eqref{eq:cell}, the positive power estimate,
and \eqref{eq:jetbound} imply
\[
 \norm{J_T(\Pi_{j+1}P-\Pi_jP)}
 \leq h\norm P_{\D}\norm{J_T}e^T
       \sum_{q=1}^d\frac{B_q+V_q}{q!}.
\]
Take the finite nonnegative sum as $L_d$, and sum the cell
differences from $a$ to $b-1$ by the triangle inequality.
The argument proves dual-norm continuity of conditional increments;
it does not assert continuity of moving-time positions in the
strong factorial source norm.
\end{proof}

\begin{proposition}\label{prop:fieldlimit}
For each fixed cylinder $P$ and every $q\geq0$, the native fields
$A_{n,q}(P)$ converge in $\mathcal F$ to a field $C_qP$.
The causal fields converge to the same limit, and the bound
\eqref{eq:Cbound} holds on the core.
\end{proposition}
\begin{proof}
Compare adjacent grids. Each coarse increment splits into two
fine increments. Differentiating the pulled-back polynomial in
either fine coordinate is the pullback of differentiating in the
coarse coordinate. Its degree stays at most $d=\deg P$, and its
factorial norm after any number of such derivatives is at most
$\norm P_{\D}$.

At the first half of a coarse cell, the strict-past projections
agree by signed convolution consistency. At the second half,
the fine conditioning cut has advanced by just one fine cell.
Apply \eqref{eq:cut} to the derivative polynomial. Dividing by
$q!$ and integrating over the unit normalized horizon gives
\begin{equation}\label{eq:adjacent}
 \norm{\widetilde A_{n+1,q}-\widetilde A_{n,q}}_{\mathcal F}
 \leq\frac{L_d}{q!}2^{-(n+1)}\norm P_{\D}.
\end{equation}
This bound is summable in $n$. Thus the causal fields are Cauchy
in the complete space $\mathcal F$ and converge. Equation
\eqref{eq:heaterror} gives the same limit for the native fields.
The norm bound follows by passing to the limit in
\eqref{eq:fieldbound}. Each native field is an adapted polynomial
step field; closedness gives $C_qP\in\mathcal A_T$.
\end{proof}

For $q>d$ all native fields are zero, hence $C_qP=0$.
If two base-grid presentations define the same cylinder, their
refinements agree on a common tail. Uniqueness of norm limits
therefore makes the limit independent of the presentation.
Finite-grid linearity and continuity of addition and scalar
multiplication make $P\mapsto C_qP$ linear. These observations
are necessary before extending the maps to completed observables.

\section{The multichannel formula and observable completion}
\begin{proposition}\label{prop:polyformula}
Every polynomial cylinder satisfies
\[
 J_T(P-\E_TP)=\sum_{q=1}^{p-1}I_q(C_qP).
\]
\end{proposition}
\begin{proof}
Refine the fixed polynomial to all sufficiently fine grids.
Choose $d'\geq\max(\deg P,p-1)$ once and pad the exact telescope
\eqref{eq:telescope} to degree $d'$ with zero coefficients.
Refinement preserves the degree bound, so this is the same finite
number of terms on every grid. The left side, after applying
$J_T$, is the same dual element on each grid.

For each $q$, \eqref{eq:jetbound} gives a uniform bound on the
paired native coefficients. Proposition~\ref{prop:fieldlimit}
provides the $L^1$ convergence for the low powers.
Apply \eqref{eq:lowl} to $q<p$, the exact critical cancellation
\eqref{eq:criticall} to $q=p$, and \eqref{eq:highl} to $q>p$.
Pass through the one fixed finite sum. The surviving terms are
exactly those in the displayed identity.
\end{proof}

\begin{proof}[Proof of Theorem~\ref{thm:main}]
The degree-independent bound \eqref{eq:Cbound} on the dense core
extends each linear coefficient map uniquely to a bounded linear
map $C_q:\D\to\mathcal F$ with the same bound. For a completed
observable choose polynomials $P_k\to u$ in $\D$.
Then $C_qP_k\to C_qu$. Closedness of $\mathcal A_T$ preserves
adapted membership. The centered map
$u\mapsto J_T(u-\E_T(u)1)$ is bounded linear, as is each $I_qC_q$.
Thus Proposition~\ref{prop:polyformula} passes to the limit and
proves \eqref{eq:main} for every $u\in\D$.

For the simultaneous actual signed approximation, assemble the
$p-1$ native fields and the native surviving power sum into one
finite product Banach space. For each fixed $P_k$, the preceding
arguments show convergence of that package as the mesh tends
to zero. Its limiting package is the collection
$(C_1P_k,\ldots,C_{p-1}P_k,J_T(P_k-\E_TP_k))$,
which tends to the analogous package for $u$ by continuity.
Choose a mesh index $n_k$ at least $k$ and at least the base-grid
index of $P_k$, large enough that the package error is less than
$1/(k+1)$. The triangle inequality gives simultaneous convergence
to the package for $u$. Every approximating sum comes from the
actual finite signed polynomial laws by \eqref{eq:actualpair}.
\end{proof}

The diagonal is a proved selection; it is not every arbitrary
joint degree/mesh limit. The convergence constants for one fixed
polynomial may depend on its degree, as $L_d$ shows. The extension
to $\D$ instead uses the degree-independent coefficient bound.
Those two kinds of bounds have distinct roles in the proof.

\section{Uniqueness and the biharmonic case}
The lower derivative orders cannot all be replaced by a coefficient
depending only on the $p$th derivative. The terminal cylinder
$P=Y_1$ has $D_s^pP=0$ for every $s$, and $\E_TP=0$.
Nevertheless
\[
 J_T(P)(Y_1^{p-1})=\E_T(Y_1^p)=cp!T\neq0.
\]
Thus any proposed reconstruction which vanishes whenever the
$p$th derivative vanishes fails on this actual signed polynomial
pairing. This elementary obstruction explains the role of lower
channels without invoking a positive pseudo-$L^2$ interpretation.

If another bounded linear coefficient map agrees with $C_q$ on
all canonical polynomial limits, it equals $C_q$ by density.
This proves uniqueness of the coefficient extension.
It does not assert that an arbitrary channel representation has
unique integrands. Let
\[
 S:(\mathcal F)^{p-1}\to\D^*,\qquad
 S(A_1,\ldots,A_{p-1})=\sum_{q=1}^{p-1}I_qA_q.
\]
Then $S(A)=S(B)$ if and only if $A-B\in\ker S$.
No theorem here sets that kernel equal to zero.

For $p=4$ the three surviving channels are
\begin{align*}
 (I_1A)(v)&=-4T\int_0^1 A(s)(D_s^3v)\,ds,\\
 (I_2A)(v)&=-12T\int_0^1 A(s)(D_s^2v)\,ds,\\
 (I_3A)(v)&=-24T\int_0^1 A(s)(D_sv)\,ds.
\end{align*}
The native coefficient fields retain factors $1$, $1/2!$, and
$1/3!$ respectively. The fourth channel is the drift
$I_4=-24TQ_0$ and cancels against the fourth moment centering.
This example keeps the signs, factorials, and complementary
test derivative orders visible.

For $p=2$ the same construction has one surviving channel
$I_1=2TQ_1$, with the kernel normalization of generator
$\partial_x^2$. It is compatible with that Brownian variance
convention, but the theorem's proof has not relied on positivity
or on importing a probability-space martingale representation.


\section{Why the reconstruction stays in the test dual}\label{sec:signed-obstructions}
The topology and channel orders in Theorem~\ref{thm:main} are essential
parts of its meaning. We give actual signed-law tests of stronger
interpretations. Except for the final terminal pairing, these statements
concern the native finite signed increment measures rather than the
factorial analytic completion. No positive sampling law is inserted.

\subsection{A finite-variation path-law obstruction}
Assume $p=2m>2$. The moment formula~\eqref{eq:moments} gives
$\int x^2k_\tau(x)dx=0$, while $\int k_\tau=1$.
If $k_\tau\geq0$ almost everywhere, the nonnegative integrable
function $x^2k_\tau(x)$ would vanish almost everywhere. Since
Lebesgue measure has no atom at zero, this would force $k_\tau=0$
almost everywhere, contradicting its unit mass. Thus the kernel is
not nonnegative. Writing $k_\tau=k_\tau^+-k_\tau^-$ shows
\[
 \int|k_\tau|=1+2\int k_\tau^->1.
\]
Every integral is finite because the kernel is Schwartz.

The inverse Fourier integral gives the exact scaling
\begin{equation}\label{eq:signed-scaling}
 k_\tau(x)=\tau^{-1/p}k_1(x\tau^{-1/p}),\qquad
 A_p:=\int_\R|k_\tau(x)|dx=\int_\R|k_1(x)|dx>1.
\end{equation}
Indeed substitute $\xi=\tau^{-1/p}\eta$ in its absolutely convergent
inverse transform and then change variables in the absolute integral.
Thus the absolute mass does not become small with the time step.
For $N$ increments on a fixed horizon $T$, the actual signed product
law $\mu_N(dx)=\prod_{j=1}^Nk_{T/N}(x_j)dx_j$ satisfies, by Tonelli,
\begin{equation}\label{eq:TVproduct}
 |\mu_N|(\R^N)=\prod_{j=1}^N\int|k_{T/N}|=A_p^N.
\end{equation}

Suppose a countably additive signed measure $\nu$ of finite total
variation on any measurable space had measurable $N$-increment
observations with these pushforward laws for every $N$. For every
measurable map $F$, $|F_\#\nu|(\text{whole space})\leq
|\nu|(\text{whole space})$. To prove this inequality, pull back any
finite measurable partition used in the variation supremum; the sum
of the absolute signed masses is bounded by the total variation of
$\nu$. Applying it to these observations would give
$A_p^N\leq|\nu|(\text{whole space})$ for all $N$, impossible.
In particular there is no such finite signed path law on $[0,T]$
realizing all the indicated uniform-grid increment distributions.
No path continuity or consistency of different observation maps was
needed for this contradiction.

The formal path-obstruction declaration uses arbitrarily many
successive increments at a fixed positive step, hence increasing
horizons. The fixed-horizon refinement conclusion just proved is
an explicit classical deduction combining the same variation
contraction with exact kernel scaling. It is not misattributed to
a different formal path declaration. For $p=2$, $A_2=1$ and this
obstruction does not apply. It concerns finite total variation;
it does not preclude generalized integration or the bounded
test-dual construction proved above. The classical obstruction is
also discussed in the primary introduction of~\cite{bcm}.

\subsection{Signed quartic cancellation versus absolute error}
Take the biharmonic order $p=4$, so $c=-1$, and put $h=T/N$.
For the actual law $\mu_N$ above define
$Q_N(x)=\sum_{j=1}^N x_j^4$. Its pointwise value is nonnegative,
but its signed mean is
\begin{equation}\label{eq:negativequartic}
 \int Q_N\,d\mu_N=N(-24h)=-24T.
\end{equation}
The one-increment moments are
$\int x^4k_h=-24h$ and $\int x^8k_h=8!h^2/2=20160h^2$.
Product integration, justified by absolute polynomial integrability,
therefore gives
\begin{align*}
 \int Q_N^2\,d\mu_N
 &=20160Nh^2+576N(N-1)h^2\\
 &=576T^2+19584T^2/N.
\end{align*}
Expanding the centered square and using~\eqref{eq:negativequartic}
yields the exact cancellation
\begin{equation}\label{eq:signedquarticerror}
 \int(Q_N+24T)^2\,d\mu_N=19584T^2/N\longrightarrow0.
\end{equation}
This is a signed integral of a square, not a norm-convergence statement.
In fact $Q_N+24T\geq24T$ everywhere. Consequently the actual absolute
error has the opposite behavior:
\begin{equation}\label{eq:absolutequarticerror}
 \int|Q_N+24T|^2\,d|\mu_N|
 \geq576T^2 A_4^N\longrightarrow\infty.
\end{equation}
Each integral is finite, since the product variation has every
polynomial moment. The absolute error even has the lower bound
$576T^2$ under the normalized variation probability
$|\mu_N|/A_4^N$. Thus neither ordinary probability convergence to
$-24T$ under those probability laws nor an $L^2$ isometry can be
read from~\eqref{eq:signedquarticerror}. The signed identity and the
divergent variation error concern precisely the same native sums.

\subsection{An $L^2$-small deterministic control with a fixed law}
The proper control scale can be tested before any completed observable
calculus. If $a$ is constant on finitely many disjoint time intervals
of lengths $h_j>0$, define its biharmonic signed integral law by the
pushforward of the product increment law under $\sum_j a_jx_j$.
Multiplying the native Fourier transforms proves
\[
 \int e^{i\xi y}\,d\mathcal L_a(y)
   =\exp\left(-\xi^4\sum_jh_ja_j^4\right).
\]
The sum is the fourth-power control clock, not its second-power norm.
On $[0,1]$ take the actual step functions
\[
 a_n(s)=(n+1)1_{(0,(n+1)^{-4}]}(s),\qquad n\geq0.
\]
They belong to both $L^2$ and $L^4$, and exactly
\[
 \|a_n\|_2=(n+1)^{-1}\longrightarrow0,\qquad
 \int_0^1 a_n^4=1.
\]
Thus $\widehat{\mathcal L_{a_n}}(1)=e^{-1}$ for every $n$, whereas
the zero control gives the Dirac law with Fourier value one.
Any scalar functional on $L^2$ agreeing with these native Fourier
values and equal to one at zero is discontinuous at zero, by this
explicit convergent sequence. A weakly continuous law extension
agreeing with the native laws is likewise impossible, since the
real and imaginary parts of $e^{iy}$ are bounded continuous tests.
This concerns a plain $L^2$ input norm for the fourth-order integral;
it does not contradict the factorial observable-domain theorem or
an extension on a stronger fourth-power control norm.

The formal source uses Fourier phases $e^{2\pi i\xi y}$, so its
unit-frequency witness is $e^{-(2\pi)^4}$ rather than $e^{-1}$.
We use the angular-frequency convention fixed at the start of this
paper. These are the same physical signed laws at correspondingly
rescaled frequencies, with the same strict failure of continuity.

\subsection{The last complementary derivative channel is essential}
Return to the analytic completion and assume $p>2$. Let $X$ be the
terminal coordinate polynomial at horizon $T$. On the polynomial
core its signed mean and pairing satisfy
\[
 \E_T(X^{p-1})=0,\qquad J_T(X^{p-1})(X)=p!cT\ne0.
\]
Both follow directly from~\eqref{eq:moments}.
For $1\leq q\leq p-2$ one has $p-q\geq2$. Since $X$ is linear,
$D_s^{p-q}X=0$ for every $s$, and the channel definition
\eqref{eq:channel} gives
\[
 (I_q A)(X)=0\qquad\text{for every admissible field }A.
\]
Every sum of those channels annihilates $X$, whereas the centered
pairing of $X^{p-1}$ does not. Hence those channels alone cannot
reconstruct that observable. In particular the first channel alone
fails. In the full $p-1$ formula, only its last channel can contribute
to this particular terminal test. This proves necessity of the last
complementary-derivative channel for the displayed reconstruction;
it does not prove an arbitrary general minimality theorem or
injectivity of channel synthesis. The original first-channel
obstruction is formalized; the statement about all earlier channels
is the same explicitly proved derivative calculation applied to
\eqref{eq:channel}.

\section{Verified source and literature boundary}
The mathematical source is the frozen Git commit
\begin{center}\small\ttfamily
bab688bb6140cd6316c80680095a914ea0d346db.
\end{center}
The namespaces in the source map share the following prefix:
\begin{center}\code{OperatorDerivativeRewrite.Foundation}\end{center}

The file \code{theorem-map.json} records exact files, declarations,
line numbers, source hashes and baseline GitHub links.

\begin{longtable}{@{}p{0.24\textwidth}p{0.70\textwidth}@{}}
\textit{Content}&\textit{Lean namespace and principal declaration}\\\hline
Domain and core&\code{HigherHeatAnalyticCylinderDomain.core_injective},
\code{core_dense}; \code{MultivariateFactorialDomain.derivative_norm}\\[3pt]
Signed product pairing&\code{HigherHeatCompletedPairing.pairing_core_complex},
\code{pairing_not_positive}\\[3pt]
Exact finite telescope&\code{HigherHeatMultitimeJetTelescope.telescope},
\code{native_dual_identity}\\[3pt]
Power bounds&\code{HigherHeatAllPowerBounds.positive_increment_bound},
\code{higher_increment_bound};
\code{HigherHeatAllChannelRemainder.cell_remainder_bound}\\[3pt]
Centered channel limits&\code{HigherHeatCenteredPowerLimits.low_centered_converge},
\code{critical_centered_vanish}, \code{higher_centered_vanish}\\[3pt]
Native/causal correction&\code{HigherHeatCausalJetIdentification.physical_shifted_jet_eq},
\code{causal_field_approximation}\\[3pt]
Cut and refinement bounds&\code{HigherHeatConditionalCutBound.conditional_cut_bound};
\code{HigherHeatCausalFieldCauchy.causal_fields_step_bound}\\[3pt]
Field limits&\code{HigherHeatMultitimePredictableFields.field_tendsto},
\code{field_bound}\\[3pt]
Polynomial reconstruction&\code{HigherHeatMultitimeReconstruction.polynomial_reconstruction}\\[3pt]
Observable extension&\code{HigherHeatMultitimeFieldMaps.completed_bound},
\code{completed_mem}; \code{HigherHeatCompletedClarkOcone.clark_ocone}\\[3pt]
Uniqueness&\code{HigherHeatCompletedClarkOcone.coefficient_unique},
\code{synthesis_uniqueness}\\[3pt]
Actual joint diagonal&\code{HigherHeatSignedDiagonalRealization.completed_signed_realization}\\
\end{longtable}

The baseline receipt records a complete successful root build and
a complete 54,943-declaration standard-axiom audit. Publication work
changes no Lean source and reuses that verified baseline. The
human proofs above reorganize the actual theorem mechanism; the
notation $B_r,L_d$ corresponds to recursive power bounds and finite
degree-dependent constants in the formal sources. It does not
claim a fresh Lean theorem merely because a constant has a shorter
name in this exposition.

The primary comparison in~\cite{bcm}, especially its discussion
leading to Theorem 18, already retains lower power increments
in a generalized It\^o calculus for complex random-walk
approximations. Its primary introduction also explicitly discusses failure of ordinary signed path-measure realization; our finite-variation obstruction is classical context, with all actual kernel/scaling details proved here. Its expected chain-rule formula is distinct from
the present completed real signed-cylinder test-dual identity.
That distinction is a comparison of the stated constructions,
not a proof of priority for the present one.
The preliminary literature check identifies earlier signed heat
calculus and alternative complex-walk approaches. It has not
established priority for this particular analytic completion or
the multitime reconstruction theorem. Hochberg's paper was
identified bibliographically; its complete text was not retrieved
in this pass, so no uninspected theorem-number comparison is made.
The cited time-change paper is a primary contemporary reference
for pseudoprocesses, not a premise extending the present result
to signed random clocks. A publication decision requires a
targeted comparison of domains and reconstruction topologies.

\begin{thebibliography}{9}
\bibitem{hochberg} K. J. Hochberg,
A signed measure on path space related to Wiener measure,
\emph{Ann. Probab.} \textbf{6} (1978), 433--458.
\bibitem{bcm} S. Bonaccorsi, C. Calcaterra and S. Mazzucchi,
An It\^o calculus for a class of limit processes arising from random
walks on the complex plane, arXiv:1603.05437 (2016).
\url{https://arxiv.org/abs/1603.05437}.
\bibitem{co} F. Cinque and E. Orsingher,
Higher-order fractional equations and related time-changed
pseudo-processes, arXiv:2402.13691v2 (2024).
\url{https://arxiv.org/abs/2402.13691}.
\end{thebibliography}
\end{document}
